\documentclass[10pt, a4paper]{article}

% ---- Packages ----
\usepackage[T1]{fontenc}
\usepackage[utf8]{inputenc}
\usepackage{lmodern}
\usepackage[margin=2.5cm]{geometry}
\usepackage{amsmath, amssymb}
\usepackage{graphicx}
\usepackage{booktabs}
\usepackage{hyperref}
\usepackage{natbib}
\usepackage{microtype}
\usepackage{titlesec}
\usepackage{parskip}
\usepackage{enumitem}
\usepackage{xcolor}
\usepackage{float}

% ---- Title ----
\title{%
  \textbf{Image2Biomass Prediction with CEMS, CHARMS, and DA-Fusion}\\[0.5em]
  \large INF367A -- CSIRO Kaggle Competition
}

\author{
  Maya Evensen \and Ragnhild Klette \and Victor Slorer
}

\date{\today}

% ---- Document ----
\begin{document}

\maketitle

%% ----------------------------------------------------------------
\begin{abstract}
We evaluate three recently published methods for small-data multi-target
regression on the CSIRO Image2Biomass Kaggle competition, where the task is to
predict five continuous biomass targets from 357 pasture photographs. Our
baseline freezes a DINOv2 ViT-S/14 backbone and trains a lightweight MLP head
with a per-target weighted Smooth-L1 loss. On top of this baseline, each team
member adapts one method targeting a different hypothesised bottleneck:
\textbf{DA-Fusion} \citep{trabucco2024} generates diffusion-based synthetic
training images; \textbf{CEMS} augments training along the local tangent space
of the joint feature-target manifold; and \textbf{CHARMS}
aligns image features with auxiliary tabular attributes available at training
time. We analyse which methods compose natively on a shared backbone and build a
combined pipeline integrating all three methods.
Individually, all three methods perform near the baseline; CEMS narrowly edges
ahead on the private leaderboard ($+0.002$ $R^{2}$) while the other two fall
slightly short. The combined pipeline underperforms all individual methods.
We discuss why each method's core assumptions interact poorly with continuous
targets, strong pretrained features, and a training set of only a few hundred
images.
\end{abstract}

%% ----------------------------------------------------------------
\section{Introduction}

Estimating pasture biomass from imagery supports grazing management and yield
forecasting, but is characterised by small labelled datasets, fine visual
distinctions between species, and continuous multi-target labels rather than
discrete classes---properties that limit what standard vision pipelines deliver
out of the box.

Self-supervised vision transformers such as DINOv2 \citep{oquab2023} offer a
natural starting point: pretrained features transfer well with limited data, and
freezing the backbone mitigates overfitting on small training sets. A frozen
DINOv2 ViT-S/14 with a lightweight MLP head forms our baseline. On top of this
baseline, each group member contributed one recently published method addressing
a different hypothesised bottleneck: DA-Fusion \citep{trabucco2024} targets data
scarcity via diffusion-based synthetic augmentation; CEMS targets geometric
structure in the joint feature-target manifold; CHARMS \citep{jiang2024} targets
the use of auxiliary tabular signal available at training but not at test time.

Our contributions are:
\begin{enumerate}[label=(\arabic*)]
  \item implementations and evaluations of each method under a unified baseline;
  \item an integration analysis identifying which methods compose;
  \item a combined pipeline integrating all three methods on a shared backbone; and
  \item a discussion of where each method's assumptions hold or break in the
        continuous-target small-data regime.
\end{enumerate}

%% ----------------------------------------------------------------
\section{Problem and Data}

The CSIRO Biomass Kaggle competition requires predicting five continuous biomass
targets from a single RGB pasture photograph: dry weights of green biomass, dead
biomass, clover, GDM, and total biomass. The metric is a weighted global
$R^{2}$, computed by flattening all (image, target) pairs into one vector with
per-target weights $0.1 / 0.1 / 0.1 / 0.2 / 0.5$; the \emph{Total} target
therefore carries half the weight.

The training set contains 357 unique pasture photographs at $2000 \times 1000$
resolution. The small size, visual similarity across species, and continuous
targets make this a challenging small-data regression problem. We use
\texttt{GroupKFold} cross-validation with five folds, grouping by source so all
orientation-augmented copies fall in the same fold; this prevents leakage
introduced by flip-based augmentation. Predictions are clipped at zero since
biomass cannot be negative.

Each training image is expanded to four orientations (identity, horizontal flip,
vertical flip, 180\textdegree), yielding 1428 training rows. This augmentation
is shared across all methods.

%% ----------------------------------------------------------------
\section{Baseline}
\label{sec:baseline}

Our baseline is a frozen self-supervised feature extractor with a lightweight
trainable head, chosen to make the most of a small training set without
overfitting. Images are resized to $504 \times 252$, normalised with ImageNet
statistics, and passed through a frozen DINOv2 ViT-S/14 backbone
\citep{oquab2023}. We use the $384$-d \texttt{[CLS]} token as the image
representation; the positional embeddings are interpolated to the non-square
input resolution.

The head, \textit{BiomassModel}, is a small MLP with an encoder
($384 \to 256 \to 64$) followed by a regression head ($64 \to 32 \to 5$). Each
linear layer is followed by GELU activation and dropout of $0.3$. Only the head
is trained, yielding roughly $117$k trainable parameters.

We train with a per-target weighted Smooth-L1 loss using the competition weights
(\{Green, Dead, Clover\}: $0.1$, GDM: $0.2$, Total: $0.5$), optimised with AdamW
at learning rate $3 \times 10^{-4}$, weight decay $10^{-3}$, batch size $64$,
for $80$ epochs. Learning rate follows a cosine schedule with
$\eta_{\min} = \eta_0 / 100$. Best checkpoints per fold are selected by
validation weighted $R^{2}$. All three method sections below modify this
baseline in one component at a time; other components remain unchanged unless
explicitly stated.

%% ----------------------------------------------------------------
\section{Methods}

\subsection{CEMS}
\label{sec:cems}

\paragraph{Motivation.}
With only 357 labelled images, the training distribution is sparse and standard
ERM trains on a fixed set of feature-target pairs. The manifold hypothesis
suggests that DINOv2 features and continuous biomass targets lie near a
low-dimensional manifold; sampling \emph{along} that manifold should yield more
meaningful synthetic points than naive interpolation. CEMS~\citep{kaufman2025}
operationalises this with a second-order local approximation that accounts for
the curvature missed by first-order methods.

\paragraph{Adaptation.}
We operate in the \emph{ambient} joint space $\mathbf{z}{=}[\mathbf{x},\mathbf{y}]$,
where $\mathbf{x}{\in}\mathbb{R}^{384}$ is the frozen DINOv2 CLS token and
$\mathbf{y}{\in}\mathbb{R}^{5}$ are the biomass targets, giving a 389-d joint
space. Because targets span very different scales (0--186\,g), we MinMax-scale
$\mathbf{y}$ to $[0,1]$ before forming the joint space and inverse-transform
synthetic labels back to raw scale for the loss.

\paragraph{Pipeline.}
For each minibatch of size $B{=}64$, CEMS centres the batch, computes a local
SVD basis $B_u$, then estimates the gradient $\nabla g$ and Hessian $\mathbf{H}$
of the normal-space embedding via ridge regression on projected neighbour
differences.  A perturbation $\boldsymbol{\eta}{\sim}\mathcal{N}(\mathbf{0},\sigma\mathbf{I}_d)$
is drawn in the $d$-dimensional tangent space and mapped back via
\begin{equation}
  \tilde{\mathbf{z}} = B_u\!\left[\boldsymbol{\eta},\;
    \boldsymbol{\eta}^{\!\top}\!\nabla g
    +\tfrac{1}{2}\boldsymbol{\eta}^{\!\top}\!\mathbf{H}\,\boldsymbol{\eta}
  \right] + \boldsymbol{\mu}_{N_z}.
\end{equation}
The intrinsic dimension $d$ is re-estimated per batch with
TwoNN~\citep{facco2017}; in practice $d{=}2$ throughout training.  We use fixed
$\sigma{=}1$ and train on CEMS-synthesised samples only (\texttt{mix\_real=False}).

\paragraph{Results.}
Table~\ref{tab:cems_results} reports validation and Kaggle performance.

\begin{table}[h]
  \centering
  \caption{Validation and Kaggle test weighted $R^2$. Higher is better.}
  \label{tab:cems_results}
  \begin{tabular}{lccc}
    \toprule
    \textbf{Model} & \textbf{Val $R^2$} & \textbf{Kaggle public} & \textbf{Kaggle private} \\
    \midrule
    Baseline & 0.7554 & 0.5533 & 0.5090 \\
    CEMS     & 0.7668 & 0.5497 & 0.5106 \\
    $\Delta$ & $+0.0114$ & $-0.0036$ & $+0.0016$ \\
    \bottomrule
  \end{tabular}
\end{table}

\paragraph{Discussion.}
CEMS is the only individual method that does not degrade private leaderboard
performance, gaining $+0.002$ $R^{2}$ over the baseline, though this margin is
within run-to-run noise. The estimated intrinsic dimension of the joint space
is $d{=}~2.19$, far below the 389-d ambient dimension, indicating that the frozen
DINOv2 backbone already captures most meaningful variation. Sampling within this
two-dimensional manifold introduces only mild diversity, limiting the benefit
of the second-order curvature correction.

\subsection{CHARMS}
\label{sec:charms}

\paragraph{Motivation.}
The CSIRO Biomass dataset includes expert-measured attributes available at
training time---pasture height, GreenSeeker NDVI, geographic state, and
dominant species---but not at inference. This matches the setting CHARMS
\citep{jiang2024} is designed for: leveraging tabular side-channel information
during training to improve an image-only model at test time.

\paragraph{Adaptation.}
The original method operates on CNN spatial feature maps; here we use the
384-dimensional CLS token from a frozen DINOv2 ViT-S/14, treating each
dimension as a channel. $k$-Means clustering ($k = 40$) groups dimensions
before the OT cost matrix is computed. Four attributes are used: two numerical
(\texttt{Pre\_GSHH\_NDVI}, \texttt{Height\_Ave\_cm}) and two categorical
(\texttt{State}, \texttt{Species\_first}), giving
$\mathbf{T} \in \mathbb{R}^{4 \times 384}$.

\paragraph{Pipeline.}
An FT-Transformer \citep{gorishniy2021} encodes tabular attributes into
per-attribute embeddings and produces an auxiliary biomass prediction. An
Optimal Transport step aligns image dimensions with attributes via inter-sample
cosine similarity, yielding transfer matrix $\mathbf{T}$. Li2t heads then
predict each attribute from $\mathbf{T}$-weighted image features. The total
loss is:
\begin{equation}
    \mathcal{L}
    =
    \mathcal{L}_{\text{image}}
    +
    \lambda_{\text{tab}}\,\mathcal{L}_{\text{tabular}}
    +
    \lambda_{\text{li2t}}\,\mathcal{L}_{\text{i2t}}
\end{equation}
with $\lambda_{\text{tab}} = 0.5$, $\lambda_{\text{li2t}} = 0.1$, and
$\mathbf{T}$ recomputed every five epochs.

\paragraph{Results.}
Table~\ref{tab:charms_results} and Figure~\ref{fig:charms_curves} report
validation and Kaggle performance over 80 epochs.

\begin{table}[h]
    \centering
    \caption{Validation and Kaggle test weighted $R^2$. Higher is better.}
    \label{tab:charms_results}
    \begin{tabular}{lccc}
        \toprule
        \textbf{Model}
            & \textbf{Val $R^2$}
            & \textbf{Kaggle public}
            & \textbf{Kaggle private} \\
        \midrule
        Baseline & 0.7935 & 0.5533 & 0.5090 \\
        CHARMS   & 0.7774 & 0.5435 & 0.5080 \\
        $\Delta$ & $-0.0161$ & $-0.0098$ & $-0.0010$ \\
        \bottomrule
    \end{tabular}
\end{table}

\begin{figure}[h]
    \centering
    \includegraphics[width=0.62\linewidth]{charms_curves.png}
    \caption{Validation curves over 80 epochs. CHARMS peaks at epoch 74
             ($R^2 = 0.7774$) versus epoch 62 for the baseline ($R^2 = 0.7697$),
             with similar convergence in both $R^2$ and loss.}
    \label{fig:charms_curves}
\end{figure}

\paragraph{Discussion.}
CHARMS does not improve over the baseline. The CLS token is lower-dimensional
than the CNN feature maps the method was designed for, limiting the quality of
the OT alignment. The dataset is also far smaller than those in the original
paper, making similarity estimates noisy. Finally, the regression setting is a
mismatch: the original paper targets classification, where discrete labels
provide stronger supervision for the auxiliary Li2t losses.


\subsection{DA-Fusion}
\label{sec:dafusion}

\paragraph{Motivation.}
DA-Fusion \citep{trabucco2024} expands small training sets by generating
synthetic images with Stable Diffusion, using textual inversion \citep{gal2023}
to learn a token per fine-grained class. It was proposed for image
classification, where each synthetic inherits its source image's class label.
We adapt it to our 5-target biomass regression setting, where labels are
continuous dry-weight vectors.

\paragraph{Adaptation.}
Textual inversion assumes discrete categories, so we use pasture species as the
grouping: one \texttt{<species\_token>} per species with at least 10 images, and
a generic prompt otherwise. The label-copy rule is stronger for regression than
for classification, since continuous targets are sensitive to visual changes
that a class label would absorb. We mitigate this by keeping img2img strength
low (0.2--0.6, jittered) so synthetics remain visually close to their source,
and apply a per-fold leakage filter excluding synthetics whose source image is
in the validation split.

\begin{figure}[h]
    \centering
    \includegraphics[width=0.8\textwidth]{example_da_fused_images.png}
    \caption{One real image (left) and three DA-Fusion synthetics.}
    \label{fig:da-fusion-examples}
\end{figure}

\paragraph{Pipeline.}
Textual inversion and image synthesis are run offline, producing 1071 synthetic
images paired with their source image's target vector. At training time,
synthetics are injected into the DINOv2\,+\,MLP baseline; all other components
(backbone, head, loss, optimiser, schedule) are unchanged.

\paragraph{Results.}
On the Kaggle leaderboard, DA-Fusion with the DINOv2 backbone scored
$0.531$ public / $0.498$ private weighted $R^{2}$, compared to the baseline's
$0.532$ / $0.509$. A ResNet-50 variant, tried as a backbone ablation, scored
$0.483$ / $0.425$. To disentangle the effect from classical pixel-space
augmentation, we ran a local fold-0 ablation against a RandAugment baseline
under the identical DINOv2 + MLP pipeline: RandAugment lifts validation
$R^{2}$ by $+0.11$, whereas DA-Fusion lifts it by only $+0.07$ and has
strictly worse per-target RMSE on Green, Clover, and Total. Both augmented
runs also converge substantially faster than the plain baseline, reaching
near-peak validation $R^{2}$ within $\sim\!15$ epochs versus $\sim\!40$ for
the baseline. The Kaggle and local trends agree: DA-Fusion helps less than
simple classical augmentation.

\paragraph{Discussion.}
Three factors likely contribute.
\textbf{(1)} The label-copy assumption is weaker in regression than in
classification: a Stable Diffusion edit that would leave a class label intact
can shift continuous biomass targets, injecting label noise.
\textbf{(2)} The DINOv2 backbone is pretrained at a scale that may already
saturate the signal extractable from 357 training images, leaving little room
for \emph{added} data to help; RandAugment still wins because it regularises
the head via feature-level perturbation rather than adding new samples.
\textbf{(3)} Synthetics outnumber real images after injection and may dilute
real-image gradients, an effect RandAugment avoids since it augments the real
images in place.
Continuous-target tasks may need a label-aware generation step (e.g.\
conditioning on the target vector) to avoid the label-noise issue.


%% ----------------------------------------------------------------
\section{Combined Pipeline}
\label{sec:combined}

The assignment asks us to combine our individual methods into a single system.
Before implementation, we analysed which of the three methods compose natively
on our shared DINOv2\,+\,MLP baseline.

\paragraph{Compatibility analysis.}
DA-Fusion and CEMS are both training-time modifications to a frozen DINOv2
backbone and compose cleanly: DA-Fusion expands the training set, and CEMS
then builds its $k$-NN graph on the expanded $[\mathbf{X}, \mathbf{Y}]$ space.
The only ordering constraint is that synthetics must be injected \emph{before}
the $k$-NN graph is built, and the target scaler must be fitted on the
combined training labels.

CHARMS is also compatible: treating the 384 CLS dimensions as channels,
as described in Section~\ref{sec:charms}, allows the OT alignment to operate
on the same frozen DINOv2 features used by the other two methods.

\paragraph{Full CHARMS integration.}
The combined pipeline runs the complete CHARMS stack on the 384-d CLS token.
An FT-Transformer encodes four tabular attributes (NDVI, pasture height,
state, species) and produces an auxiliary biomass prediction; K-Means
($k{=}40$) clusters image dimensions by co-activation; Earth Mover Distance
aligns channel- and attribute-similarity matrices to yield
$\mathbf{T}\!\in\!\mathbb{R}^{4\times 384}$, recomputed every five epochs.
CHARMS losses are applied only to real samples, as synthetics carry no
tabular data.

\paragraph{Combined training loop.}
For each \texttt{GroupKFold} split: (i) filter synthetics to the training
split; (ii) extract DINOv2 features and tabular attributes for real training
rows; (iii) fit the target scaler on the combined labels; (iv) cluster 384
image dimensions via K-Means and initialise $\mathbf{T}$ uniformly;
(v) train with the CEMS anchor loop---CHARMS losses computed on the
real-sample subset of each anchor batch, $\mathbf{T}$ refreshed via EMD every
five epochs. Checkpoint selection by validation weighted $R^{2}$.

\subsection{Results}
\label{sec:combined-results}

Table~\ref{tab:combined-results} reports the combined pipeline's performance
on the Kaggle leaderboard alongside the baseline and each individual method
for reference.

\begin{table}[h]
\centering
\small
\begin{tabular}{lcc}
\toprule
Pipeline & Public $R^{2}_{w}$ & Private $R^{2}_{w}$ \\
\midrule
Baseline (DINOv2\,+\,MLP)            & \textbf{0.553}& 0.509\\
DA-Fusion                            & 0.531 & 0.498 \\
CEMS                                 & 0.549& \textbf{0.510}\\
CHARMS                               & 0.543& 0.508\\
\midrule
Combined (DA-Fusion\,+\,CEMS\,+\,CHARMS) & 0.549& 0.485\\
\bottomrule
\end{tabular}
\caption{Weighted $R^{2}$ on the Kaggle leaderboard. The combined pipeline
does not recover baseline performance.}
\label{tab:combined-results}
\end{table}

The combined pipeline scores $0.549 / 0.485$ public/private weighted $R^{2}$---matching
CEMS on the public leaderboard but falling below all individual methods on the private
score. Combining the three modifications does not recover baseline performance.

\subsection{Discussion}
\label{sec:combined-discussion}

The combined pipeline confirms the pattern visible in the individual results:
each method's modification reduces performance, and stacking them compounds
rather than cancels the effect. This is consistent with each method
addressing a \emph{hypothesised} bottleneck that turns out not to be the real
one in this setting---data scarcity for DA-Fusion, latent-space geometry for
CEMS, and cross-modal supervision for CHARMS. When the actual constraint is
elsewhere (plausibly the quality of the 5-target label signal given 357
images and a frozen backbone that already saturates available representation
capacity), adding components that each introduce their own noise or
regularisation pressure makes the system worse, not better.

\paragraph{Limitations.}
Our conclusions come with caveats. Each method is evaluated with the
hyperparameters reported in its source paper plus minimal tuning; a more
thorough search might reduce the gap to baseline. Kaggle leaderboard scores
reflect a single training run per method rather than a seed-averaged mean,
so differences of $\pm 0.01$ $R^{2}$ are within run-to-run noise.

%% ----------------------------------------------------------------
\section{Conclusion}
We adapted three recently published methods to a 5-target pasture-biomass
regression task with 357 training images and a frozen DINOv2 backbone.
Individually, all three methods perform near the baseline; CEMS narrowly edges
it out on the private leaderboard ($+0.002$ $R^{2}$) while DA-Fusion and CHARMS
fall slightly short. The combined pipeline underperforms all individual methods. We attribute
this to mismatches between each method's target domain (classification,
geometric regularity, or cross-modal supervision) and the actual constraints
of our setting, and suggest that continuous-target small-data regression
problems may benefit more from label-aware generation and head-level
regularisation than from the forms of augmentation tested here.

%% ----------------------------------------------------------------
\bibliographystyle{plainnat}
\begin{thebibliography}{9}

\bibitem[Facco et~al.(2017)]{facco2017}
Facco, E., d'Errico, M., Rodriguez, A., and Laio, A. (2017).
\newblock Estimating the intrinsic dimension of datasets by a minimal
  neighborhood information.
\newblock \emph{Scientific Reports}, 7(1):12140.

\bibitem[Gal et~al.(2023)]{gal2023}
Gal, R., Alaluf, Y., Atzmon, Y., Patashnik, O., Bermano, A.~H., Chechik, G.,
and Cohen-Or, D. (2023).
\newblock An image is worth one word: Personalizing text-to-image generation
  using textual inversion.
\newblock In \emph{Proceedings of the Eleventh International Conference on
  Learning Representations (ICLR)}.

\bibitem[Kaufman \& Azencot(2025)]{kaufman2025}
Kaufman, I. and Azencot, O. (2025).
\newblock Curvature enhanced data augmentation for regression.
\newblock In \emph{Proceedings of the 42nd International Conference on Machine
  Learning (ICML)}, PMLR 267.

\bibitem[Gorishniy et~al.(2021)]{gorishniy2021}
Gorishniy, Y., Rubachev, I., Khrulkov, V., and Babenko, A. (2021).
\newblock Revisiting deep learning models for tabular data.
\newblock In \emph{Advances in Neural Information Processing Systems
  (NeurIPS)}, volume~34, pages 18932--18943.

\bibitem[Jiang et~al.(2024)]{jiang2024}
Jiang, J., Ye, H., Wang, L., Yang, Y., Jiang, Y., and Zhan, D. (2024).
\newblock Tabular insights, visual impacts: Transferring expertise from tables
  to images.
\newblock In \emph{Proceedings of the 41st International Conference on Machine
  Learning}, PMLR 235.

\bibitem[Oquab et~al.(2023)]{oquab2023}
Oquab, M., Darcet, T., Moutakanni, T., et~al. (2023).
\newblock {DINOv2}: Learning robust visual features without supervision.
\newblock \emph{arXiv preprint arXiv:2304.07193}.

\bibitem[Trabucco et~al.(2024)]{trabucco2024}
Trabucco, B., Doherty, K., Gurinas, M., and Salakhutdinov, R. (2024).
\newblock Effective data augmentation with diffusion models.
\newblock In \emph{Proceedings of the Twelfth International Conference on
  Learning Representations (ICLR)}.

\end{thebibliography}

\end{document}